\begin{proposition}[Proved anticanonical support comparison]\label{cdp:st:C.2.VII.1}
\Uses{the hypotheses, notation and definitions stated here; \cite[Theorem~0.2, Proposition~2.29 and Theorem~2.38]{Dloussky2006}; \cite[\S1, pp.~284--285, and \S3, p.~296]{DOT2003}; \cite[Corollary~2.25 and Proposition~2.27]{Dloussky2006}; \cite[Introduction, p.~283]{DOT2003}; \cite[Lemma~2.22 and the following paragraph; Proposition~2.27]{Dloussky2006}; \cref{cdp:st:C.3.8}}



Let $Z$ be a smooth compact surface of class VII, not
necessarily minimal, and let $D>0$ be an effective integral
divisor such that $K_Z+D$ is torsion in $\Pic(Z)$.
Write $R=D_{\mathrm{red}}$, and let $r(R)$ be its number
of rational irreducible components. Then
\[
 e(R)=r(R),\qquad r(Z)\le b_2(Z),
\tag{C.VII.1}\label{cdp:eq:C.VII.1}
\]
where $r(Z)$ counts all rational curves of $Z$.
Every connected component of $R$ has arithmetic genus one.
In particular a connected component containing an elliptic
curve consists of that single elliptic curve with trees of
rational curves attached; it cannot contain a second
elliptic curve or a path joining two distinct points of the
first one.
\end{proposition}
\begin{proof}
\noindent\textbf{Step 1.} We specify the established surface results used in this
comparison. For a minimal class VII surface with positive
second Betti number, existence of a numerically
anticanonical divisor implies existence of a global
spherical shell by
\cite[Theorem~0.2, Proposition~2.29 and Theorem~2.38]{Dloussky2006}.
The reduced support is one rational cycle with attached
trees, one or two rational cycles, or, in the parabolic
Inoue case, a disjoint rational cycle and a smooth elliptic
curve. The intermediate graph and its actual curve
configuration are described in
\cite[\S1, pp.~284--285, and \S3, p.~296]{DOT2003}.
The fact that no connected piece of the anticanonical
support is an isolated tree is
\cite[Corollary~2.25 and Proposition~2.27]{Dloussky2006};
Lemma~2.22 there shows that a curve meeting that support
belongs to it. These statements apply to our divisor with
index one, since its twisting line is torsion and therefore
flat. They give, on the minimal model, arithmetic genus
one for each connected component and Euler number equal
to the number of its rational components. The rational
curve count is at most $b_2$, as stated and used in
\cite[Introduction, p.~283]{DOT2003}.
These are consequences of the proved theorem about
anticanonical divisors; no classification assertion about
class VII surfaces without such a divisor is used.

\noindent\textbf{Step 2.} There is a coefficient issue in the parabolic Inoue case
which the torsion hypothesis resolves explicitly. Write
$E$ for its elliptic curve and $C$ for its reduced rational
cycle. The exact canonical relation is
$K_{Z_0}+E+C=0$; the curves $E$ and $C$ are disjoint,
$E^2=-b_2(Z_0)$, and the intersection matrix of the
components of $C$ has kernel generated by their sum
\cite[Lemma~2.22 and the following paragraph; Proposition~2.27]{Dloussky2006}.
Consequently an effective integral divisor with
$c_1(K_{Z_0}+D_0)_{\mathbb R}=0$ has the form
$D_0=E+kC$, $k\ge0$: intersection with $E$ fixes its
coefficient at one, and the cycle intersection matrix
makes all cycle coefficients equal. Our stronger
torsion relation then gives
$K_{Z_0}+D_0=(k-1)C$ torsion. Choose $n>0$ trivializing $\OO((k-1)C)^{\otimes n}$.
A trivialization compared with the canonical meromorphic
section gives a meromorphic function with divisor
$n(k-1)C$. For $k>1$ it has no poles and a nonzero
effective zero divisor, impossible on the compact
connected surface. For $k=0$ its reciprocal has that
property. Hence $k=1$. In particular both connected
pieces of the stated reduced support really occur.
This deduction uses the torsion relation, not merely
the vanishing of its real first Chern class.

\noindent\textbf{Step 3.} For a minimal model with $b_2=0$, the nonzero effective
divisor excludes the curve-free Inoue surfaces. The
Hopf alternative has only smooth elliptic curves, which
are pairwise disjoint. This follows as well from the
explicit contraction description used in \ref{cdp:st:C.3.8} and its
finite unramified quotient: the preimages of two distinct
curves are disjoint elliptic curves, and the quotient
is locally biholomorphic. All assertions are then
immediate on that model.

\noindent\textbf{Step 4.} We verify the transfer through the actual blowdowns.
Let $\sigma:Z\to Z_0$ be a sequence of point blowups
from a minimal model. A torsion line is represented by
a finite unitary character of the fundamental group,
and a point blowup induces an isomorphism on fundamental
groups. Thus the torsion line $K_Z+D$ descends to $Z_0$.
Pushing forward a meromorphic canonical section with pole
divisor $D$, or its torsion-trivializing power, gives an
effective divisor $D_0$ with $K_{Z_0}+D_0$ torsion.
It is nonzero, since a torsion canonical bundle would
give Kodaira dimension zero. At each individual blowup
the canonical divisor formula and the pushforward section
give the exact divisor equality
\[
 D_{\mathrm{new}}=\sigma^*D_{\mathrm{old}}-E.
\tag{C.VII.2}\label{cdp:eq:C.VII.2}
\]
The coefficient of $E$ is
$\operatorname{mult}_pD_{\mathrm{old}}-1$.
Effectivity forces the centre to lie on $D_{\mathrm{old}}$.
If this coefficient is positive, the reduced support
is its entire reduced inverse image: a smooth point
acquires a rational leaf, or an ordinary intersection
is subdivided by the exceptional curve. These operations
preserve arithmetic genus and increase both the Euler
number and the rational-component count by one.
If the coefficient is zero, the centre is a smooth point
of a multiplicity-one component and the new support is
its strict transform, isomorphic to the old support.
Thus all the stated support properties transfer in the
required direction. Finally each blowup increases $b_2$
by one and adds precisely its exceptional rational
curve; every other rational curve is the strict transform
of one on the preceding surface. The curve-count
inequality in \eqref{cdp:eq:C.VII.1} follows.
\end{proof}







